\subsection*{全体概要：この章で学ぶこと}
\addcontentsline{toc}{subsection}{全体概要：この章で学ぶこと}
この章は、ソフトウェアを構成する品目を識別し、変更を統制し、状態を記録し、監査し、正しい内容をビルドしてリリースするための知識を扱う。ソフトウェア構成管理は、版管理ツールの操作だけではなく、ソフトウェアの完全性、正確性、追跡可能性、説明責任を保つための管理体系である。

\begin{notebox}{到達目標}この章を読み終えると、構成品目、基準版、変更要求、構成状態記録、構成監査、リリース管理の役割を説明できるようになる。また、なぜ構成管理が品質保証、セキュリティ、保守、規制対応に関わるのかを説明できるようになる。\end{notebox}
\begin{itemize}
\item 「何を管理対象にするか」を構成品目として定められる。
\item 構成品目を変更する前に、変更要求、影響分析、承認、実施、検証の流れを説明できる。
\item 基準版が、変更を止めるものではなく、変更を管理可能にする基点であることを説明できる。
\item 状態記録・報告が、証跡、品質保証、セキュリティ、規制対応に使われることを説明できる。
\item ビルド、部品表、リリースノート、版記述文書がリリース管理で果たす役割を説明できる。
\end{itemize}
\par\smallskip
\begin{center}
\centering
\IfFileExists{assets/generated_figures/ch08/fig-ch08-01.png}{\includegraphics[width=0.96\linewidth]{assets/generated_figures/ch08/fig-ch08-01.png}}{\fbox{\parbox[c][48mm][c]{0.9\linewidth}{\centering 画像生成待ち\\第8章の全体地図}}}
\par\smallskip
\refstepcounter{figure}\label{fig:ch08-01}図 \thefigure: 第8章の全体地図
\end{center}
図 \thefigure は、第8章の全体地図を視覚的に整理したものである。
\par\smallskip
